\documentclass[11pt,a4paper]{article}
\usepackage[T1]{fontenc}
\usepackage[utf8]{inputenc}
\usepackage{lmodern,microtype}
\usepackage[a4paper,margin=26mm,headheight=15pt]{geometry}
\usepackage{amsmath,amssymb,amsthm,mathtools,bm,mathrsfs}
\usepackage{booktabs,array,longtable,enumitem,ragged2e,xurl}
\usepackage[dvipsnames]{xcolor}
\usepackage{fancyhdr,needspace}
\usepackage[unicode,colorlinks=true,linkcolor=MidnightBlue,citecolor=MidnightBlue,urlcolor=MidnightBlue]{hyperref}
\usepackage[nameinlink,noabbrev]{cleveref}
\usepackage{listings}
\hypersetup{pdftitle={q-Binomial Coefficient Calculus},pdfauthor={Research extension},pdfsubject={Gaussian binomials, Bell polynomials, interpolation, inversion, and root-of-unity derivatives}}
\numberwithin{equation}{section}
\newtheorem{theorem}{Theorem}[section]
\newtheorem{proposition}[theorem]{Proposition}
\newtheorem{lemma}[theorem]{Lemma}
\newtheorem{corollary}[theorem]{Corollary}
\theoremstyle{definition}
\newtheorem{definition}[theorem]{Definition}
\newtheorem{example}[theorem]{Example}
\theoremstyle{remark}
\newtheorem{remark}[theorem]{Remark}
\newcommand{\Q}{\mathbb Q}
\newcommand{\Z}{\mathbb Z}
\newcommand{\N}{\mathbb N}
\newcommand{\C}{\mathbb C}
\newcommand{\qb}[2]{\genfrac{[}{]}{0pt}{}{#1}{#2}_{q}}
\newcommand{\qbb}[3]{\genfrac{[}{]}{0pt}{}{#1}{#2}_{#3}}
\newcommand{\qf}[1]{[#1]_{q}!}
\newcommand{\qi}[1]{[#1]_{q}}
\newcommand{\poch}[3]{(#1;#2)_{#3}}
\newcommand{\coef}[2]{[#1^{#2}]}
\newcommand{\Ecal}{\mathcal E}
\newcommand{\Pcal}{\mathcal P}
\newcommand{\Bhat}{\widehat B}
\newcommand{\Bq}{\mathcal B^{(q)}}
\newcommand{\Res}{\mathop{\mathrm{Res}}}
\newcommand{\ee}{\mathrm e}
\newcommand{\ii}{\mathrm i}
\newcommand{\stir}[2]{\genfrac{\{} {\}}{0pt}{}{#1}{#2}}
\newcommand{\sone}[2]{s(#1,#2)}
\DeclareMathOperator{\diag}{diag}
\DeclareMathOperator{\ord}{ord}
\DeclareMathOperator{\Id}{Id}
\pagestyle{fancy}
\fancyhf{}
\fancyhead[L]{\small\itshape q-Binomial Coefficient Calculus}
\fancyhead[R]{\small\thepage}
\renewcommand{\headrulewidth}{0.35pt}
\setlength{\parskip}{4pt plus 1pt}
\setlength{\parindent}{1.2em}
\setlength{\emergencystretch}{2em}
\allowdisplaybreaks[2]
\setcounter{tocdepth}{2}
\lstset{basicstyle=\ttfamily\small,breaklines=true,columns=fullflexible,frame=single,showstringspaces=false}

\begin{document}
\begin{titlepage}
\thispagestyle{empty}
\vspace*{20mm}
\begin{center}
{\small\scshape An extension to Combinatorial Coefficient Calculus}\\[12mm]
{\Huge\bfseries q-Binomial\\[3mm] Coefficient Calculus}\\[8mm]
{\Large Gaussian weights, Bell expansions,\\[2mm]
interpolation, reversion, and root-of-unity jets}\\[12mm]
{\large September 4, 2026}
\end{center}
\vspace{10mm}
\begin{abstract}
Gaussian binomial coefficients fit into coefficient calculus in several different,
noninterchangeable ways: as elementary and complete symmetric functions on geometric
alphabets, as convolution structure constants, as interpolation matrices on
$q$-integer nodes, and as polynomials whose local expansions are governed by
Bernoulli, Stirling, and Bell numbers. This article develops those connections with
explicit hypotheses and proofs. A finite partition-sum kernel evaluates arbitrary
powers of products of $q$-shifted factorials, their compositions, and their
Lagrange--B\"urmann inverses. Further results include weighted $q$-binomial inversion,
two Newton interpolation formulas, a normalization-specific $q$-Stirling calculus,
ordinary composition in $q$-divided-power coordinates, coupled two-variable
inversion, exact geometric extrapolation, and formulas for every derivative of a
Gaussian polynomial at every root of unity. Classical identities are distinguished
from the organizing synthesis. In particular, a change of factorial normalization
is not confused with a Jackson chain rule, and distinct $q$-Catalan constructions
are not identified. A reproducible exact-arithmetic verification program accompanies
the source and rendered article; its finite tests supplement, rather than replace,
the proofs.
\end{abstract}
\vfill
\noindent\textbf{Scope.} A self-contained mathematical extension, not a modification
of the repository and not a claim of new Lean formalization. All coefficient
formulas have explicitly finite expansions at each prescribed degree.
\end{titlepage}
\tableofcontents
\newpage

\section{Purpose, source alignment, and algebraic conventions}\label{sec:scope}

\subsection{What is being extended}
The canonical \emph{Combinatorial Coefficient Calculus} manuscript in the ProveIt
repository \cite{CCC} connects polynomial basis changes, Bell-polynomial
composition, binomial inversion, and Lagrange--B\"urmann/Good inversion. The present
article extends these mechanisms to Gaussian coefficients. It uses the canonical
\TeX{} source consulted on September 4, 2026, rather than assuming that the
repository's retained PDF has current source parity; the provenance document
explicitly warns against that assumption \cite{provenance}.

The point is not to replace every $n!$ by $\qf n$. There are three separate
operations. One may weight an ordinary combinatorial object by a statistic, change
the basis used to record an ordinary formal series, or replace the derivative by
a difference operator. Each produces useful $q$-formulas, but they need not produce
the same formulas. Keeping these operations separate is essential for correct
composition and inversion.

The classical $q$-binomial and terminating summation identities are standard
\cite{DLMF17,DLMF175,DLMF176}. They are proved here because their proofs reveal the
coefficient mechanisms used later. The principal contribution of this extension
is the explicit integration of these identities with the source's coefficient
calculus, together with worked finite-sum consequences. No historical priority
claim is made for identities obtained by combining standard tools.

\subsection{Coefficient rings and boundary values}
Throughout, $\N=\{0,1,2,\ldots\}$. Unless stated otherwise, $q$ is an
indeterminate and computations take place in a
commutative $\Q(q)$-algebra $R$. Additional commuting parameters may be adjoined.
The main series variable is $z$, so $R[[z]]$ means formal power series in $z$,
not in $q$. An identity in $\Z[q]$ can subsequently be specialized at any complex
$q$, including a root of unity. A rational expression can be specialized only
where it is defined, or after its removable singularities have been eliminated.

Write
\begin{equation}\label{eq:definitions}
 \qi n=1+q+\cdots+q^{n-1},\qquad
 \qf n=\prod_{j=1}^n\qi j,\qquad
 D_n=\poch q q n=(1-q)^n\qf n,
\end{equation}
with $\qi0=0$, $\qf0=D_0=1$. The shifted factorial and Gaussian coefficient are
\begin{equation}\label{eq:gauss-def}
 \poch a q N=\prod_{j=0}^{N-1}(1-aq^j),\qquad
 \qb nk=\frac{D_n}{D_kD_{n-k}}\quad(0\le k\le n).
\end{equation}
For $n\ge0$, put $\qb nk=0$ when $k<0$ or $k>n$. Negative upper entries, when used,
will receive a separate definition. Empty products are $1$, empty sums are $0$,
and $0^0=1$ in finite coefficient formulas. Thus every zeroth coefficient below
is unambiguous. A coefficient of negative degree is zero.

A series $A$ with $A(0)=1$ has the formal powers
\begin{equation}\label{eq:formal-power}
 A^\alpha=\exp(\alpha\log A),\qquad
 \log(1+u)=\sum_{j\ge1}\frac{(-1)^{j-1}}j u^j.
\end{equation}
These require a $\Q$-algebra, but no analytic branch. Analytic interpretations are
stated separately in \cref{sec:analytic}.

\subsection{The ordinary Bell kernel retained throughout}
For $n\ge0$ let
\[
 \Pcal(n)=\{(m_1,\ldots,m_n)\in\N^n:\ \sum_{j=1}^n jm_j=n\}.
\]
For $n=0$ this consists of the empty tuple. Define
\begin{equation}\label{eq:E-kernel}
 \Ecal_n(L)=\sum_{\bm m\in\Pcal(n)}
       \prod_{j=1}^n\frac{(L_j/j)^{m_j}}{m_j!}.
\end{equation}
The complete exponential Bell polynomial $B_n$ is normalized by
\begin{equation}\label{eq:Bell-def}
 \exp\left(\sum_{j\ge1}x_j\frac{z^j}{j!}\right)
  =\sum_{n\ge0}B_n(x_1,\ldots,x_n)\frac{z^n}{n!}.
\end{equation}
Consequently
\begin{equation}\label{eq:E-Bell}
 \Ecal_n(L)=\frac1{n!}B_n(0!L_1,1!L_2,\ldots,(n-1)!L_n),\qquad
 \exp\left(\sum_{j\ge1}\frac{L_j}jz^j\right)
       =\sum_{n\ge0}\Ecal_n(L)z^n.
\end{equation}
Both statements follow by multiplying the individual exponential series. In
particular, \eqref{eq:E-kernel} is a definition by a finite sum, not by a recurrence.
The first values provide a useful normalization check:
\begin{align}\label{eq:E-small}
 \Ecal_0&=1,&\Ecal_1&=L_1,&
 \Ecal_2&=\frac{L_1^2+L_2}{2},\\
 \Ecal_3&=\frac{L_1^3+3L_1L_2+2L_3}{6},&
 \Ecal_4&=\frac{L_1^4+6L_1^2L_2+3L_2^2+8L_1L_3+6L_4}{24}.
 \nonumber
\end{align}
The $q$-dependence will usually enter through the arguments $L_j$, not through a
new Bell-polynomial definition.

\section{Gaussian coefficients as coefficient extractors}\label{sec:gaussian}

\subsection{The finite theorem, polynomiality, and symmetry}
\begin{theorem}[Finite $q$-binomial theorem]\label{thm:finite}
For $N\ge0$,
\begin{equation}\label{eq:finite}
 \prod_{j=0}^{N-1}(1+zq^j)
 =\sum_{k=0}^{N}q^{\binom k2}\qb Nk z^k.
\end{equation}
Moreover,
\begin{equation}\label{eq:subset}
 \qb Nk=\sum_{0\le i_1<\cdots<i_k<N}
          q^{i_1+\cdots+i_k-\binom k2}\in\Z_{\ge0}[q].
\end{equation}
The $k=0$ sum has one empty term.
\end{theorem}
\begin{proof}
Expanding the product gives the subset sum before division by $q^{\binom k2}$.
Let $A_{N,k}$ denote that normalized sum. Separating subsets that contain $N-1$
from those that do not gives
\[
 A_{N,k}=A_{N-1,k}+q^{N-k}A_{N-1,k-1}.
\]
The quotient in \eqref{eq:gauss-def} satisfies the same identity: after using the
common denominator $D_kD_{N-k}$, the numerator is
$D_{N-1}((1-q^{N-k})+q^{N-k}(1-q^k))=D_N$.
The boundary values agree, so induction proves the identification. Finally,
$i_j\ge j-1$, hence every exponent in \eqref{eq:subset} is nonnegative.
\end{proof}

\begin{corollary}[Pascal, symmetry, reciprocity, and classical limit]\label{cor:pascal}
For the indicated nonnegative indices,
\begin{align}
 \qb Nk&=\qb{N-1}k+q^{N-k}\qb{N-1}{k-1}
       =q^k\qb{N-1}k+\qb{N-1}{k-1},\label{eq:pascal}\\
 \qb Nk&=\qb N{N-k},\qquad
 \qbb Nk{q^{-1}}=q^{-k(N-k)}\qb Nk,\label{eq:reciprocity}\\
 \qbb Nk1&=\binom Nk,\qquad \qbb Nk0=1\quad(0\le k\le N).
 \label{eq:specializations}
\end{align}
For $0\le k\le N$ the degree is $k(N-k)$ and the leading and constant
coefficients are $1$.
\end{corollary}
\begin{proof}
The first Pascal identity was proved above. Symmetry follows from the factorial
quotient and gives the second Pascal identity. Replacing each $1-q^{-j}$ by
$-q^{-j}(1-q^j)$ gives reciprocity. At $q=1$, the subset sum counts $k$-subsets.
Its smallest exponent, zero, occurs once, and its largest exponent,
$k(N-k)$, occurs once. This also proves the assertion at zero.
\end{proof}

Put $\lambda_j=i_j-(j-1)$. Then
$0\le\lambda_1\le\cdots\le\lambda_k\le N-k$, and the exponent in
\eqref{eq:subset} is $\sum_j\lambda_j$. Thus $\qb Nk$ enumerates partitions
inside a $k\times(N-k)$ rectangle by size; see also \cite{DLMF26}. This interpretation
will justify positivity of coefficients that later emerge from alternating sums.

\subsection{Complete symmetric functions and negative upper entries}
\begin{theorem}[Reciprocal finite product]\label{thm:reciprocal}
For $N\ge1$,
\begin{equation}\label{eq:reciprocal-product}
 \frac1{\poch z q N}=\sum_{r\ge0}\qb{N+r-1}r z^r.
\end{equation}
Equivalently,
\begin{equation}\label{eq:complete-symmetric}
 \qb{N+r-1}r
 =\sum_{\substack{a_0+\cdots+a_{N-1}=r\\a_j\ge0}}
       q^{\sum_{j=0}^{N-1}j a_j}
 =\sum_{0\le i_1\le\cdots\le i_r<N}q^{i_1+\cdots+i_r}.
\end{equation}
\end{theorem}
\begin{proof}
The first finite sum comes from expanding every geometric series. A weakly
increasing $r$-tuple is carried to a strictly increasing tuple by
$i_j\mapsto i_j+j-1$. Its entries lie in $\{0,\ldots,N+r-2\}$ and its weight
increases by $\binom r2$. Applying \eqref{eq:subset} proves the result.
\end{proof}
For $N=0$, the reciprocal product is $1$; this case is not encoded by a Gaussian
coefficient with upper entry $r-1$.

For $M\ge1$ and $r\ge0$ define a negative upper entry by the finite product
\[
 \qb{-M}r:=\prod_{j=0}^{r-1}\frac{1-q^{-M-j}}{1-q^{j+1}}.
\]
Factoring each numerator gives the explicit conversion
\begin{equation}\label{eq:negative-upper}
 \qb{-M}r=(-1)^r q^{-Mr-\binom r2}\qb{M+r-1}r.
\end{equation}
Only the upper index has been extended here. This avoids the competing
conventions that arise when both indices are negative; a broader treatment is
available in \cite{FormichellaStraub}.

\subsection{Gaussian multinomials and a finite-field interpretation}
For nonnegative $k_1+\cdots+k_s=N$ set
\[
 \qbb{N}{k_1,\ldots,k_s}{q}=\frac{D_N}{D_{k_1}\cdots D_{k_s}}.
\]
It is a polynomial because it factors as
$\qb N{k_1}\qb{N-k_1}{k_2}\cdots\qb{k_s}{k_s}$.
When $q=Q$ is a prime power, $\qbb NkQ$ counts $k$-dimensional subspaces of
$\mathbb F_Q^N$: count ordered linearly independent $k$-tuples and divide by the
number of ordered bases of one such subspace, obtaining
\[
 \frac{\prod_{j=0}^{k-1}(Q^N-Q^j)}{\prod_{j=0}^{k-1}(Q^k-Q^j)}
 =\qbb NkQ.
\]
Successively choosing subspaces of prescribed quotient dimensions gives the
multinomial interpretation as a flag count. This is a distinct interpretation
from weighting subsets; both evaluate the same polynomial.

\section{Convolution and terminating summation identities}\label{sec:convolutions}

\subsection{Two-block and multiblock Vandermonde identities}
\begin{theorem}[$q$-Vandermonde]\label{thm:vandermonde}
For $M,N,r\ge0$,
\begin{equation}\label{eq:vandermonde}
 \qb{M+N}r
 =\sum_k q^{(M-k)(r-k)}\qb Mk\qb N{r-k}.
\end{equation}
The sum is over $\max(0,r-N)\le k\le\min(M,r)$.
More generally, if $N_1+\cdots+N_s=N$, then
\begin{equation}\label{eq:multivandermonde}
 \qb Nr=
 \sum_{\substack{k_1+\cdots+k_s=r\\0\le k_i\le N_i}}
 q^{\sum_{i<j}(N_i-k_i)k_j}\prod_{i=1}^s\qb{N_i}{k_i}.
\end{equation}
\end{theorem}
\begin{proof}
Split $\poch{-z}q{M+N}$ as $\poch{-z}qM\poch{-zq^M}qN$ and use
\eqref{eq:finite}. The coefficient with $k$ choices from the first block is
$q^{\binom k2+\binom{r-k}2+M(r-k)}\qb Mk\qb N{r-k}$.
Divide by $q^{\binom r2}$ to obtain \eqref{eq:vandermonde}.
For several blocks, their initial exponents are
$0,N_1,N_1+N_2,\ldots$. The same calculation subtracts
$\sum_{i<j}k_i k_j$ from $\sum_{i<j}N_i k_j$, proving
\eqref{eq:multivandermonde}.
\end{proof}

\begin{corollary}[Three useful companions]\label{cor:companions}
For $M,N\ge1$ and $r\ge0$,
\begin{align}
 \qb{M+N+r-1}r
 &=\sum_{k=0}^r q^{M(r-k)}\qb{M+k-1}k\qb{N+r-k-1}{r-k},
 \label{eq:negative-vander}\\
 \delta_{r0}
 &=\sum_{k=0}^{\min(N,r)}(-1)^kq^{\binom k2}
     \qb Nk\qb{N+r-k-1}{r-k}.\label{eq:orth-product}
\end{align}
For $0\le k\le n$,
\begin{equation}\label{eq:hockey}
 \sum_{j=k}^n q^{j-k}\qb jk=\qb{n+1}{k+1}.
\end{equation}
\end{corollary}
\begin{proof}
The first formula compares coefficients in
$1/\poch zq{M+N}=1/(\poch zqM\poch{zq^M}qN)$.
The second compares coefficients in
$\poch zqN/\poch zqN=1$.
For the third, rearrange the first Pascal identity as
$q^{j-k}\qb jk=\qb{j+1}{k+1}-\qb j{k+1}$ and telescope.
\end{proof}

\subsection{The infinite theorem with a precise formal meaning}
\begin{theorem}[The $q$-binomial series]\label{thm:infinite}
In $\Q(q,a)[[z]]$ define
\begin{equation}\label{eq:R-log}
 R_a(z)=\exp\left(\sum_{j\ge1}\frac{1-a^j}{j(1-q^j)}z^j\right).
\end{equation}
Then
\begin{equation}\label{eq:qbinomial-series}
 R_a(z)=\sum_{n\ge0}\frac{\poch a q n}{D_n}z^n
       =\frac{\poch{az}q\infty}{\poch zq\infty}.
\end{equation}
The last expression denotes either its coefficientwise $q$-adic interpretation
or, for $|q|<1$ and $|z|<1$, the convergent analytic product. It is not a limit of
finite products in the discrete-coefficient $z$-adic topology alone.
\end{theorem}
\begin{proof}
Subtract the logarithm at $qz$ from the logarithm at $z$ in \eqref{eq:R-log}.
Exponentiation gives
\[
 (1-z)R_a(z)=(1-az)R_a(qz).
\]
If $c_n=[z^n]R_a$, this implies
$(1-q^n)c_n=(1-aq^{n-1})c_{n-1}$ for $n\ge1$, with $c_0=1$.
Multiplying these finitely many ratios yields $c_n=\poch a q n/D_n$.
Expanding logarithms of the infinite products gives \eqref{eq:R-log}; this
calculation is valid coefficientwise $q$-adically and in the indicated analytic
domain. For finite products the coefficient of $z$ contains
$1+q+\cdots+q^{N-1}$, which does not eventually stabilize in $\Q(q)$.
This proves the topology warning as well.
\end{proof}

\begin{proposition}[Cocycle and shifted-factorial convolution]
\label{prop:cocycle}
One has $R_a(z)R_b(az)=R_{ab}(z)$ and therefore
\begin{equation}\label{eq:poch-conv}
 \sum_{k=0}^n\qb nk\poch a q k\poch b q{n-k}a^{n-k}
       =\poch{ab}q n.
\end{equation}
\end{proposition}
\begin{proof}
In the logarithm the summands combine as
$(1-a^j)+a^j(1-b^j)=1-(ab)^j$. Extracting $z^n$ and multiplying by $D_n$
proves the convolution.
\end{proof}

\begin{theorem}[The two terminating $q$-Chu--Vandermonde sums]\label{thm:chu}
For $n\ge0$, as rational identities in generic parameters,
\begin{align}
 \sum_{k=0}^n\frac{\poch{q^{-n}}qk\poch a q k}
                         {D_k\poch c q k}\,q^k
 &=a^n\frac{\poch{c/a}q n}{\poch c q n},\label{eq:chu1}\\
 \sum_{k=0}^n\frac{\poch{q^{-n}}qk\poch a q k}
                         {D_k\poch c q k}
                    \left(\frac{cq^n}{a}\right)^k
 &=\frac{\poch{c/a}q n}{\poch c q n}.\label{eq:chu2}
\end{align}
Specializations require the displayed denominators to be nonzero or a separately
justified cancellation. In \eqref{eq:chu1}, the right side can be read as
$\prod_{j=0}^{n-1}(a-cq^j)/\poch c q n$, avoiding an artificial singularity at $a=0$.
\end{theorem}
\begin{proof}
In \eqref{eq:poch-conv} replace its parameters by $A,b$, interchange $k$ and
$n-k$, and divide by $\poch A q n$. Use
\[
 \qb nk=\frac{\poch{q^{-n}}qk}{D_k}(-1)^kq^{nk-\binom k2},\qquad
 \frac{\poch A q{n-k}}{\poch A q n}
  =\frac{(-A)^{-k}q^{-nk+k(k+1)/2}}
         {\poch{q^{1-n}/A}qk}.
\]
The result is the left side of \eqref{eq:chu1} with $b=a$ and
$c=q^{1-n}/A$. Reversing the factors in the quotient on the right gives
$\poch{Aa}q n/\poch A q n=a^n\poch{c/a}q n/\poch c q n$.
For \eqref{eq:chu2}, replace $(q,a,c)$ by $(q^{-1},a^{-1},c^{-1})$ in
\eqref{eq:chu1} and use
\[
 \poch{x}{q^{-1}}k=(-x)^kq^{-\binom k2}\poch{x^{-1}}qk.
\]
The powers remaining in the summand are $(cq^n/a)^k$; those on the right cancel.
\end{proof}

\section{The master product-to-Bell theorem}\label{sec:master}

\subsection{Geometric power sums}
For a finite $N\ge0$ define
\begin{equation}\label{eq:S}
 S_N(t)=\sum_{j=0}^{N-1}t^j.
\end{equation}
This polynomial, rather than $(1-t^N)/(1-t)$, is the preferred expression at
$t=1$. When $N=\infty$, write $S_\infty(t)=(1-t)^{-1}$ with the formal or
analytic interpretation of \cref{thm:infinite}.

\begin{theorem}[Ramified shifted-factorial coefficient calculus]\label{thm:master}
Let $d_r,h_r$ be positive integers, $N_r\in\N\cup\{\infty\}$, and let
$a_r,\alpha_r$ be commuting parameters, $1\le r\le s$. Set
\begin{equation}\label{eq:master-P}
 P(z)=\prod_{r=1}^s\poch{a_r z^{d_r}}{q^{h_r}}{N_r}^{\alpha_r},
\end{equation}
and define, by finite sums for each $j\ge1$,
\begin{equation}\label{eq:master-L}
 L_j=-\sum_{\substack{1\le r\le s\\d_r\mid j}}
       d_r\alpha_r a_r^{j/d_r}
             S_{N_r}\!\left(q^{h_rj/d_r}\right).
\end{equation}
Then for every $n\ge0$ and every scalar $\beta$,
\begin{equation}\label{eq:master-coef}
 [z^n]P(z)^\beta
 =\Ecal_n(\beta L)
 =\sum_{\bm m\in\Pcal(n)}\prod_{j=1}^n
                   \frac{(\beta L_j/j)^{m_j}}{m_j!}.
\end{equation}
Here infinite factors are defined by their logarithms over $\Q(q)$; finite
factors need no such qualification.
\end{theorem}
\begin{proof}
For one finite factor,
\[
 \log\poch{az^d}{q^h}N
  =-\sum_{\ell\ge1}\frac{a^\ell}{\ell}
               S_N(q^{h\ell})z^{d\ell}.
\]
The same formula defines an infinite factor with $S_\infty$. In degree $j$, a
factor contributes only if $d\mid j$; writing $\ell=j/d$ makes its contribution
$-d a^{j/d}S_N(q^{hj/d})/j$. Summing the exponents $\alpha_r$ yields
$\log P=\sum_{j\ge1}L_jz^j/j$. Equation \eqref{eq:E-Bell} proves
\eqref{eq:master-coef}. All coefficient extractions use only $L_1,\ldots,L_n$.
\end{proof}

\begin{corollary}[Gaussian coefficients as explicit Bell sums]\label{cor:gauss-bell}
For $N\ge1$ and $k\ge0$,
\begin{align}
 q^{\binom k2}\qb Nk
 &=\sum_{\bm m\in\Pcal(k)}\prod_{j=1}^k
    \frac1{m_j!}\left(\frac{(-1)^{j-1}S_N(q^j)}j\right)^{m_j},
 \label{eq:elementary-Bell}\\
 \qb{N+k-1}k
 &=\sum_{\bm m\in\Pcal(k)}\prod_{j=1}^k
    \frac1{m_j!}\left(\frac{S_N(q^j)}j\right)^{m_j}.
 \label{eq:complete-Bell}
\end{align}
The first expression is zero for $k>N$.
\end{corollary}
\begin{proof}
Apply \cref{thm:master} to $\poch{-z}qN$ and $1/\poch zqN$, respectively,
and compare with \cref{thm:finite,thm:reciprocal}.
\end{proof}
These are the geometric-alphabet specializations of the source's
symmetric-function/Bell correspondence. No $q$-Bell polynomial is needed.
At $q=1$ they become finite Bell-sum formulas for $\binom Nk$ and
$\binom{N+k-1}k$.

\subsection{Composition, arbitrary powers, and parameter derivatives}
\begin{corollary}[Composition with a unit product]\label{cor:compose-P}
If $H(u)=\sum_{k\ge0}h_ku^k$ and $P$ is as in \cref{thm:master}, then
\begin{equation}\label{eq:compose-P}
 [z^n]H(P(z)-1)
 =\sum_{k=0}^n h_k\sum_{j=0}^k(-1)^{k-j}\binom kj\Ecal_n(jL).
\end{equation}
If $M(\bm m)=\sum_jm_j$, then
\begin{equation}\label{eq:beta-derivative}
 [z^n]P(z)^\beta(\log P(z))^r
 =\sum_{\bm m\in\Pcal(n)}
   (M(\bm m))_{\underline r}\,\beta^{M(\bm m)-r}
   \prod_{j=1}^n\frac{(L_j/j)^{m_j}}{m_j!},
\end{equation}
where terms with $M(\bm m)<r$ are zero, so no negative power of $\beta$ is used.
\end{corollary}
\begin{proof}
Expand $(P-1)^k=\sum_{j=0}^k(-1)^{k-j}\binom kjP^j$.
Since $P-1$ has positive $z$-order, only $k\le n$ can contribute.
For the second formula differentiate \eqref{eq:master-coef} $r$ times in $\beta$;
its degree in $\beta$ is at most $n$.
\end{proof}

For example, put
\[
 P(z)=\frac{\poch{-z^2}{q^2}M^\gamma}{\poch{az}qN^\alpha}.
\]
The logarithmic inputs are explicitly
\begin{equation}\label{eq:mixed-example-L}
 L_j=\alpha a^jS_N(q^j)
   -\mathbf1_{2\mid j}\,2\gamma(-1)^{j/2}S_M(q^j).
\end{equation}
Thus all powers, logarithmic factors, and compositions of this mixed even/ordinary
product are evaluated by one finite kernel. This example also makes clear why the
ramification factor $d_r$ in \eqref{eq:master-L} cannot be omitted.

\section{Extracting coefficients in the base parameter}\label{sec:qcoeff}
The preceding sections treated $q$ as a parameter. We now extract powers of $q$
itself. These are different coefficient questions.

\begin{theorem}[Divisor-sum formula for Gaussian coefficients]\label{thm:qcoeff}
Fix $0\le k\le n$ and put
\begin{equation}\label{eq:divisor-L}
 A_m(n,k)=\sum_{\substack{1\le d\le k\\d\mid m}}d
          -\sum_{\substack{n-k+1\le d\le n\\d\mid m}}d.
\end{equation}
Then, for every $r\ge0$,
\begin{equation}\label{eq:qcoeff-Bell}
 [q^r]\qb nk
 =\sum_{\bm m\in\Pcal(r)}\prod_{j=1}^r
                 \frac{(A_j(n,k)/j)^{m_j}}{m_j!}.
\end{equation}
In particular the expression vanishes for $r>k(n-k)$.
\end{theorem}
\begin{proof}
As a unit of $\Q[[q]]$,
\[
 \qb nk=\frac{\prod_{d=n-k+1}^n(1-q^d)}{\prod_{d=1}^k(1-q^d)}.
\]
The coefficient of $q^m$ in its logarithm is $A_m(n,k)/m$, since
$[q^m]\log(1-q^d)=-d/m$ when $d\mid m$ and is zero otherwise.
Apply \eqref{eq:E-Bell}, now with $q$ in place of $z$. The degree bound was
proved in \cref{cor:pascal}.
\end{proof}

\begin{proposition}[Bounded-part inclusion--exclusion]\label{prop:restricted}
Let
\[
 p_{\le k}(r)=\#\{(a_1,\ldots,a_k)\in\N^k:\ \sum_{j=1}^k ja_j=r\},
\]
with value zero for $r<0$, and $p_{\le0}(r)=\delta_{r0}$. Then
\begin{equation}\label{eq:qcoeff-inclusion}
 [q^r]\qb nk=
 \sum_{S\subseteq\{1,\ldots,k\}}(-1)^{|S|}
 p_{\le k}\!\left(r-\sum_{i\in S}(n-k+i)\right).
\end{equation}
\end{proposition}
\begin{proof}
Expand the numerator of the product in the preceding proof over subsets $S$.
The reciprocal denominator has coefficient $p_{\le k}(r)$ by the geometric
series expansion. Their Cauchy product gives the formula.
\end{proof}

For example,
\begin{equation}\label{eq:G42}
 \qb42=1+q+2q^2+q^3+q^4.
\end{equation}
Both \eqref{eq:qcoeff-Bell} and \eqref{eq:qcoeff-inclusion} recover these
coefficients without a Gaussian recurrence. Their cancellations above degree
four are exact identities. The subset formula, not the alternating presentation,
explains why the surviving coefficients are nonnegative integers.

\section{Gaussian convolution and the correct inverse transform}\label{sec:transforms}

\subsection{Divided-power coordinates and two exponentials}
Define
\begin{equation}\label{eq:qexp}
 e_q(z)=\sum_{n\ge0}\frac{z^n}{\qf n},\qquad
 E_q(z)=\sum_{n\ge0}q^{\binom n2}\frac{z^n}{\qf n}.
\end{equation}
Using \cref{thm:infinite} with $a=0$ and the finite theorem followed by its
coefficientwise limit gives
\begin{equation}\label{eq:qexp-products}
 e_q(z)=\frac1{\poch{(1-q)z}q\infty},\qquad
 E_q(z)=\poch{-(1-q)z}q\infty,
 \qquad e_q(z)E_q(-z)=1.
\end{equation}
For completeness, the coefficientwise limit used for the second equality is
$q^{\binom n2}/D_n$: in the finite theorem,
$\qb Nn=\prod_{j=0}^{n-1}(1-q^{N-j})/D_n$ tends $q$-adically to $1/D_n$.
Thus \eqref{eq:qexp-products} also holds purely formally over $\Q(q)$.
The functions in \eqref{eq:qexp} use $q$-factorials; a convention with $D_n$ in
the denominator rescales their argument by $1-q$.

For sequences $a,b$ put
\begin{equation}\label{eq:qconv}
 (a\star_q b)_n=\sum_{k=0}^n\qb nk a_kb_{n-k},\qquad
 \mathscr G_q(a;z)=\sum_{n\ge0}a_n\frac{z^n}{\qf n}.
\end{equation}
Multiplication of formal series immediately proves
\begin{equation}\label{eq:qconv-series}
 \mathscr G_q(a\star_q b;z)=\mathscr G_q(a;z)\mathscr G_q(b;z).
\end{equation}
Hence $\star_q$ is associative and commutative and has identity
$\epsilon_n=\delta_{n0}$. Alternatively, associativity follows directly from
$\qb nk\qb kj=\qb nj\qb{n-j}{k-j}$, so it is a polynomial identity and does
not depend on the invertibility of $q$-factorials after specialization.

\begin{theorem}[Weighted Gaussian inversion]\label{thm:transform}
Let $c$ be a scalar. The two transforms
\begin{align}
 b_n&=\sum_{k=0}^n\qb nk c^{n-k}a_k,\label{eq:transform-forward}\\
 a_n&=\sum_{k=0}^n(-1)^{n-k}q^{\binom{n-k}2}\qb nk c^{n-k}b_k
 \label{eq:transform-inverse}
\end{align}
are mutually inverse. These are polynomial identities, valid over every
commutative ring after assigning a value to $q$ and $c$.
\end{theorem}
\begin{proof}
Over $\Q(q)$, the forward transform multiplies the generating series by
$e_q(cz)$ and the inverse multiplies by $E_q(-cz)$. Their product is $1$.
For a denominator-free proof, the coefficient of $a_j$ after both transforms is
\[
 c^{n-j}\qb nj
 \sum_{r=0}^{n-j}(-1)^r q^{\binom r2}\qb{n-j}r
 =c^{n-j}\qb nj\poch1q{n-j}.
\]
This is $1$ when $n=j$ and $0$ otherwise. Interchanging the transforms gives the
same sum, proving both inverse compositions. Since the entire argument is a
polynomial identity with integer coefficients, arbitrary specialization follows.
\end{proof}

\subsection{Why ordinary translation fails}
Let $T_c$ denote \eqref{eq:transform-forward}. Define the homogeneous polynomial
\[
 H_r(c,d;q)=\sum_{j=0}^r\qb rj c^j d^{r-j}.
\]
Associativity of convolution gives the exact composition law
\begin{equation}\label{eq:transform-composition}
 (T_cT_d a)_n=\sum_{k=0}^n\qb nk H_{n-k}(c,d;q)a_k.
\end{equation}
For example,
\[
 H_2(c,d;q)=c^2+(1+q)cd+d^2,
\]
which is not $(c+d)^2$ for generic commuting $c,d$. Thus
$T_cT_d\ne T_{c+d}$ in general, and $T_{-c}$ is not the inverse of $T_c$.
The missing factor in the naive inverse is precisely $q^{\binom{n-k}2}$.
This contrasts with the ordinary weighted-translation theorem in the base article.
The correct replacement is either \eqref{eq:transform-composition}, or the
multiplicative cocycle of \cref{prop:cocycle} when shifted-factorial kernels are used.

\subsection{Subspace-lattice M\"obius inversion}
For a prime power $Q$, let $U\subseteq V$ be finite-dimensional vector spaces and
$r=\dim(V/U)$. The M\"obius function of their subspace interval is
\begin{equation}\label{eq:subspace-mobius}
 \mu(U,V)=(-1)^r Q^{\binom r2}.
\end{equation}
Indeed, the sum of this expression over intermediate subspaces of an
$r$-dimensional quotient is
$\sum_{j=0}^r\qbb rjQ(-1)^jQ^{\binom j2}=\poch1Qr$.
It is $1$ for $r=0$ and zero otherwise, exactly the inverse-to-zeta condition.
Grouping subspace sums by dimension produces \cref{thm:transform} with $c=1$.
Thus the extra Gaussian inverse weight has a structural explanation: the Boolean
lattice has M\"obius weight $(-1)^r$, whereas the subspace lattice has the additional
factor $Q^{\binom r2}$.

\section{Newton interpolation and a \texorpdfstring{$q$}{q}-Stirling basis change}\label{sec:newton}

\subsection{The \texorpdfstring{$q$}{q}-integer lattice: the deformation of ordinary Newton expansion}
Define the monic polynomials
\begin{equation}\label{eq:Fk}
 F_k(x)=\prod_{i=0}^{k-1}(x-\qi i),\qquad F_0(x)=1.
\end{equation}
Their normalized versions
\[
 N_k(x;q)=\frac{F_k(x)}{q^{\binom k2}\qf k}
\]
satisfy $N_k(\qi m;q)=\qb mk$ for $m\ge0$. To see this, use
$\qi m-\qi i=q^i\qi{m-i}$ and multiply over $i$.
This basis belongs to the broader setting of quantum integer-valued polynomials
\cite{HarmanHopkins}; its normalization is stated here explicitly.

\begin{theorem}[Finite $q$-Newton expansion]\label{thm:qnewton}
For every polynomial $p$ of degree at most $d$ over $R$,
\begin{align}
 p(x)&=\sum_{k=0}^d c_k N_k(x;q),\label{eq:qnewton}\\
 c_k&=\sum_{j=0}^k(-1)^{k-j}q^{\binom{k-j}2}\qb kj p(\qi j).
 \label{eq:qnewton-coef}
\end{align}
At $q=1$ this becomes
$p(x)=\sum_{k=0}^d\Delta^kp(0)\binom xk$.
\end{theorem}
\begin{proof}
The distinct degrees and unit leading coefficients of $N_0,\ldots,N_d$ make
them a basis. Evaluate its expansion at $x=\qi m$ for $0\le m\le d$:
$p(\qi m)=\sum_{k=0}^m\qb mk c_k$. Gaussian inversion gives
\eqref{eq:qnewton-coef}. At $q=1$, the factors in $N_k$ become
$(x-i)$, their denominator becomes $k!$, and the sum for $c_k$ is the ordinary
finite-difference formula. All denominators in this specialization are nonzero.
\end{proof}

\subsection{A specified pair of \texorpdfstring{$q$}{q}-Stirling arrays}
\begin{definition}\label{def:qstirling}
Define $S_q(n,k)$ and $s_q(n,k)$ by
\begin{equation}\label{eq:qstirling-def}
 x^n=\sum_{k=0}^nS_q(n,k)F_k(x),\qquad
 F_n(x)=\sum_{k=0}^n s_q(n,k)x^k.
\end{equation}
These definitions, rather than the words ``$q$-Stirling numbers'' alone, fix the
convention. Both arrays are zero outside $0\le k\le n$ and have entry $1$ at
$(0,0)$.
\end{definition}

\begin{theorem}[Finite formulas and inverse matrices]\label{thm:qstirling}
For $n,k\ge0$ with $k\le n$,
\begin{align}
 S_q(n,k)
 &=\frac1{q^{\binom k2}\qf k}
   \sum_{j=0}^k(-1)^{k-j}q^{\binom{k-j}2}\qb kj(\qi j)^n,
 \label{eq:qstirling-explicit}\\
 s_q(n,k)
 &=(-1)^{n-k}
   \sum_{\substack{I\subseteq\{0,\ldots,n-1\}\\|I|=n-k}}
         \prod_{i\in I}\qi i.\label{eq:qstirling-first}
\end{align}
Furthermore,
\begin{equation}\label{eq:qstirling-inverse}
 \sum_{j=k}^nS_q(n,j)s_q(j,k)=\delta_{nk},\qquad
 \sum_{j=k}^ns_q(n,j)S_q(j,k)=\delta_{nk}.
\end{equation}
The arrays belong to $\Z[q]$, although \eqref{eq:qstirling-explicit} is displayed
as a rational function. In fact $S_q(n,k)\in\Z_{\ge0}[q]$.
\end{theorem}
\begin{proof}
Apply \cref{thm:qnewton} to $p(x)=x^n$ to obtain the first formula. Expand the
product $F_n$ over subsets to obtain the second. Substitution of the two basis
expansions into each other gives both matrix identities.
To prove polynomiality without division, note that
$xF_k=F_{k+1}+\qi k F_k$. Starting from $x^0=F_0$ and multiplying by $x$
shows that every $S_q(n,k)$ is a sum of products of $q$-integers with nonnegative
integer coefficients. This argument also yields the check
\[
 S_q(n+1,k)=S_q(n,k-1)+\qi k S_q(n,k),
\]
but the explicit formula \eqref{eq:qstirling-explicit} does not depend on using
this recurrence to evaluate earlier rows.
\end{proof}

A second finite expression is
\begin{equation}\label{eq:qstirling-h}
 S_q(n,k)=
 \sum_{\substack{a_1+\cdots+a_k=n-k\\a_j\ge0}}
       \prod_{j=1}^k(\qi j)^{a_j},\qquad n\ge k.
\end{equation}
To prove it, the coefficient sequence on the right has generating series
$\prod_{j=1}^k(1-\qi j\,t)^{-1}$. The recurrence in the preceding proof implies,
for $G_k(t)=\sum_{n\ge k}S_q(n,k)t^{n-k}$, that
$(1-\qi k\,t)G_k(t)=G_{k-1}(t)$, with $G_0=1$; hence the products agree.
When $k=0$, \eqref{eq:qstirling-h} means $\delta_{n0}$.
All three formulas degenerate to the corresponding ordinary Stirling formulas at
$q=1$. For example, $S_q(3,2)=2+q$.

\subsection{Geometric nodes are different}
Let $P_k(x)=\prod_{i=0}^{k-1}(x-q^i)$. For $q\ne0$ with the relevant powers
$q^0,\ldots,q^d$ distinct, every polynomial $p$ of degree at most $d$ has
\begin{equation}\label{eq:geometric-newton}
 p(x)=\sum_{k=0}^d A_kP_k(x),\qquad
 A_k=\sum_{j=0}^k
       \frac{(-1)^j q^{-kj+j(j+1)/2}}{D_jD_{k-j}}p(q^j).
\end{equation}
Indeed,
$P_k(q^m)=(-1)^kq^{\binom k2}D_k\qb mk$; apply the same triangular inversion.
Equivalently, $A_k$ is the usual divided difference on the geometric nodes,
because
\[
 \prod_{\substack{0\le i\le k\\i\ne j}}(q^j-q^i)
   =(-1)^j q^{jk-j(j+1)/2}D_jD_{k-j}.
\]
This supplies an independent interpolation proof of \eqref{eq:geometric-newton}.

As $q\to1$, these nodes all coalesce at $1$ and $A_k\to p^{(k)}(1)/k!$;
thus the limiting expansion is Taylor's formula at $1$, not ordinary Newton
expansion on the integers. One algebraic proof of the limit is that the monic
bases $P_k$ tend coefficientwise to $(x-1)^k$; their finite triangular
change-of-basis matrix has diagonal $1$ and hence a continuously varying inverse.
The distinction between the two lattices is important in applications.

\section{Ordinary composition in \texorpdfstring{$q$}{q}-divided-power coordinates}\label{sec:composition}

\subsection{A normalized Bell family, with ordinary multiplicities}
The ordinary partial Bell polynomial is
\[
 \Bhat_{n,k}(a_1,a_2,\ldots)
 =\sum_{\substack{\sum_jm_j=k\\\sum_jjm_j=n}}
          \frac{k!}{\prod_jm_j!}\prod_ja_j^{m_j}.
\]
It satisfies $[z^n](\sum_{j\ge1}a_jz^j)^k=\Bhat_{n,k}(a)$.
Define $\Bhat_{0,0}=1$ and $\Bhat_{n,0}=0$ for $n>0$.

\begin{definition}\label{def:Bq}
For ordinary substitution of commuting series, define
\begin{equation}\label{eq:Bq-def}
 \Bq_{n,k}(g)
 =\frac{\qf n}{\qf k}
   \sum_{\substack{\sum_jm_j=k\\\sum_jjm_j=n}}
       \frac{k!}{\prod_jm_j!}
       \prod_j\left(\frac{g_j}{\qf j}\right)^{m_j}.
\end{equation}
Boundary values are $\Bq_{0,0}=1$ and $\Bq_{n,0}=0$ for $n>0$.
\end{definition}
This is a coordinate-normalized Bell family; it is not asserted to coincide with
any other family called ``$q$-Bell polynomials.'' Its exact relation to the base
article's ordinary Bell polynomial is
\begin{equation}\label{eq:Bq-conversion}
 \Bq_{n,k}(g)=\frac{\qf n}{\qf k}
           \Bhat_{n,k}\left(\frac{g_1}{\qf1},\frac{g_2}{\qf2},\ldots\right).
\end{equation}

\begin{theorem}[Composition rule]\label{thm:qcomposition}
If
$F(z)=\sum_{k\ge0}f_kz^k/\qf k$ and
$G(z)=\sum_{j\ge1}g_jz^j/\qf j$, then
\begin{equation}\label{eq:qcomposition}
 \qf n\,[z^n]F(G(z))=\sum_{k=0}^nf_k\Bq_{n,k}(g).
\end{equation}
At $q=1$, $\Bq_{n,k}$ becomes the exponential partial Bell polynomial $B_{n,k}$.
\end{theorem}
\begin{proof}
Expand $G^k$ using the ordinary multinomial theorem. Its multiplicity for the
profile $(m_j)$ is $k!/\prod_jm_j!$; there is no Gaussian weight because the
factors commute and ordinary substitution is being used. Multiply by
$f_k/\qf k$, sum over $k$, and multiply by $\qf n$. At $q=1$, cancellation of
$k!$ yields the usual exponential Bell profile $n!/\prod_jm_j!(j!)^{m_j}$.
\end{proof}

The first nontrivial examples expose a frequent normalization error:
\begin{align}
 \Bq_{3,2}&=\frac{2\qi3}{\qi2}g_1g_2,\label{eq:Bq32}\\
 \Bq_{4,2}&=\frac{2\qi4}{\qi2}g_1g_3
             +\frac{\qi4\qi3}{(\qi2)^2}g_2^2,\label{eq:Bq42}\\
 \Bq_{4,3}&=\frac{3\qi4}{\qi2}g_1^2g_2.
 \label{eq:Bq43}
\end{align}
Also $\Bq_{n,1}=g_n$ and $\Bq_{n,n}=g_1^n$. The rational coefficient
$2\qi3/\qi2$ is not $\qi3$. It need not be a polynomial in $q$, and it may have a
pole at a root of unity. This is expected: the coordinate system itself uses
inverses of $q$-factorials.

\subsection{Riordan matrices as a coordinate change}
\begin{theorem}[Normalized Riordan composition]\label{thm:riordan}
Let $A(0)$ and $G'(0)$ be units and $G(0)=0$. Define
\begin{equation}\label{eq:riordan-def}
 \mathcal R_q(A,G)_{n,k}=\frac{\qf n}{\qf k}[z^n]A(z)G(z)^k.
\end{equation}
Then
\begin{align}
 \mathcal R_q(A,G)\mathcal R_q(B,H)
  &=\mathcal R_q(A\cdot B\circ G,H\circ G),\label{eq:riordan-product}\\
 \mathcal R_q(A,G)^{-1}
  &=\mathcal R_q\left(\frac1{A\circ\overline G},\overline G\right),
 \label{eq:riordan-inverse}
\end{align}
where $\overline G$ is the ordinary compositional inverse.
\end{theorem}
\begin{proof}
The matrix acts on the coefficient vector of $F$ by sending its series to
$A(z)F(G(z))$. Applying this operation twice proves
\eqref{eq:riordan-product}; substituting $H=\overline G$ and
$B=1/(A\circ\overline G)$ proves \eqref{eq:riordan-inverse}.
Equivalently, if $R(A,G)_{n,k}=[z^n]A G^k$, then
$\mathcal R_q=D R D^{-1}$ for $D=\diag(\qf0,\qf1,\ldots)$.
Thus this family is a diagonal renormalization of ordinary Riordan matrices,
not an unexplained new substitution law.
\end{proof}
In particular the Gaussian transform is
$\mathcal R_q(e_q(cz),z)$ and its inverse is
$\mathcal R_q(E_q(-cz),z)$, recovering \cref{thm:transform}.

\section{Jackson calculus and the noncommuting binomial theorem}\label{sec:jackson}

Define the Jackson derivative and dilation by
\begin{equation}\label{eq:Jackson}
 D_qf(z)=\frac{f(z)-f(qz)}{(1-q)z},\qquad T_qf(z)=f(qz).
\end{equation}
The numerator has zero constant term, so the quotient is a formal series.
One has $D_qz^n=\qi n z^{n-1}$, $D_qT_q=qT_qD_q$, and
\begin{equation}\label{eq:qleibniz1}
 D_q(fg)=(D_qf)g+(T_qf)(D_qg).
\end{equation}
These follow by direct subtraction.

\begin{theorem}[Iterated Jackson product rule]\label{thm:jackson}
For $n\ge0$,
\begin{equation}\label{eq:qleibniz}
 D_q^n(fg)(z)=\sum_{k=0}^n\qb nk
       (D_q^kf)(q^{n-k}z)(D_q^{n-k}g)(z).
\end{equation}
\end{theorem}
\begin{proof}
The case $n=0$ is immediate. Differentiate each summand using
\eqref{eq:qleibniz1}. The derivative falling on its first factor contributes
$q^{n-k}(D_q^{k+1}f)(q^{n-k}z)D_q^{n-k}g$; the derivative falling on its second
factor contributes $(D_q^kf)(q^{n-k+1}z)D_q^{n-k+1}g$.
After reindexing, the coefficient of the new $k$th term is
$\qb nk+q^{n+1-k}\qb n{k-1}=\qb{n+1}k$ by Pascal's identity.
\end{proof}

At the origin, $D_q^nf(0)=\qf n[z^n]f$. Consequently
\cref{thm:qcomposition} gives an exact origin formula for
$D_q^n(F\circ G)(0)$ when $G(0)=0$. It does \emph{not} imply an ordinary-looking
Jackson chain rule away from zero. The valid first-order identity is
\begin{equation}\label{eq:qchain}
 D_q(F\circ G)(z)=D_qG(z)\,
       F[G(z),G(qz)],
\end{equation}
where $F[u,v]=(F(u)-F(v))/(u-v)$ is the divided difference. Its polynomial or
formal expansion removes the apparent denominator when $u=v$.
For $F(w)=w^2$ this is
$D_qG(z)(G(z)+G(qz))$, not generally $\qi2 G(z)D_qG(z)$.
For instance $G(z)=z^2$ makes the former second factor $(1+q^2)z^2$.

There is, however, a genuine Gaussian binomial theorem for noncommuting variables.
\begin{proposition}[Quantum-plane binomial theorem]\label{prop:quantum-plane}
In an associative algebra with a central parameter $q$ and $YX=qXY$,
\begin{equation}\label{eq:quantum-plane}
 (X+Y)^n=\sum_{k=0}^n\qb nk X^{n-k}Y^k.
\end{equation}
\end{proposition}
\begin{proof}
Right-multiply by $X+Y$ and use $Y^kX=q^kXY^k$.
The new coefficient is $q^k\qb nk+\qb n{k-1}=\qb{n+1}k$.
The initial case is $n=0$.
\end{proof}
This theorem is consistent with the failure of ordinary scalar translation in
\cref{sec:transforms}: the Gaussian cross term is produced by a commutation
relation, not by the ordinary expansion of commuting $c+d$.

\section{Lagrange inversion with shifted-factorial weights}\label{sec:lagrange}

\subsection{A fully explicit coefficient theorem}
\begin{theorem}[Product-weighted Lagrange--B\"urmann]\label{thm:lagrange}
Let $\Phi(t)$ be a product of the form \eqref{eq:master-P}, with logarithmic
inputs $L_j$ from \eqref{eq:master-L}. There is a unique $w(z)\in zR[[z]]$
satisfying $w=z\Phi(w)$. For integers $n\ge m\ge1$,
\begin{equation}\label{eq:lagrange-master}
 [z^n]w(z)^m
 =\frac mn\Ecal_{n-m}(nL)
 =\frac mn\sum_{\bm a\in\Pcal(n-m)}\prod_{j=1}^{n-m}
                      \frac{(nL_j/j)^{a_j}}{a_j!}.
\end{equation}
For $H(t)=\sum_{r\ge0}h_rt^r$ and $n\ge1$,
\begin{equation}\label{eq:lagrange-H}
 [z^n]H(w(z))=\frac1n\sum_{r=1}^n r h_r\Ecal_{n-r}(nL).
\end{equation}
The constant coefficient is $h_0$. For $n<m$ the coefficient of $w^m$ is zero.
\end{theorem}
\begin{proof}
The map $t\mapsto t/\Phi(t)=t+O(t^2)$ has a unique formal inverse, which is
$w$. Alternatively, iteration from $w=0$ improves the $z$-adic precision at every
step and proves existence and uniqueness without analytic assumptions.
For the coefficient formula, make the formal residue substitution
$z=t/\Phi(t)$:
\begin{align*}
 [z^n]w^m
 &=\Res_{t=0}t^{m-n-1}\Phi(t)^n
              \left(1-t\frac{\Phi'(t)}{\Phi(t)}\right)\,dt\\
 &=[t^{n-m}]\Phi^n-\frac1n[t^{n-m-1}](\Phi^n)'
   =\frac mn[t^{n-m}]\Phi^n.
\end{align*}
The residue change is valid for a substitution with linear coefficient $1$;
it can also be checked on each Laurent monomial, whose derivative has residue
zero except for the logarithmic differential. Apply \cref{thm:master} to the last
coefficient. Expanding $H$ and using linearity gives \eqref{eq:lagrange-H}.
\end{proof}
This is ordinary Lagrange inversion over a coefficient field containing $q$.
Neither the prefactor $m/n$ nor the derivative in Lagrange--B\"urmann is replaced
by a $q$-integer or a Jackson derivative.

\subsection{A finite Gaussian convolution form}
For $p,N\ge0$ define
\begin{equation}\label{eq:Cpn}
 C_r^{(p,N)}(q)=
 \sum_{\substack{k_1+\cdots+k_p=r\\0\le k_i\le N}}
       q^{\sum_{i=1}^p\binom{k_i}2}\prod_{i=1}^p\qb N{k_i}
 =[t^r]\poch{-t}qN^p.
\end{equation}
For $p=0$ or $N=0$ this is $\delta_{r0}$. The equality follows by applying the
finite theorem to each of the $p$ factors; this definition requires no recurrence.

\begin{corollary}[A slot-weighted Fuss family]\label{cor:fuss}
If $w=z\poch{-w}qN$, then
\begin{equation}\label{eq:slot-fuss}
 [z^n]w^m=\frac mn C_{n-m}^{(n,N)}(q)\qquad(n\ge m\ge1).
\end{equation}
At $q=1$ this is $\frac mn\binom{Nn}{n-m}$. In particular, for $N\ge1$,
\begin{equation}\label{eq:fuss-limit}
 [z^n]w\big|_{q=1}
 =\frac1n\binom{Nn}{n-1}
 =\frac1{(N-1)n+1}\binom{Nn}n.
\end{equation}
\end{corollary}
\begin{proof}
Use \cref{thm:lagrange} in coefficient-extraction form and then
\eqref{eq:Cpn}. At $q=1$, $\poch{-t}1N=(1+t)^N$; the last equality is an
ordinary factorial cancellation.
\end{proof}
Although \eqref{eq:slot-fuss} contains $m/n$, the case $m=1$ has coefficients in
$\Z_{\ge0}[q]$. A direct interpretation proves this: a rooted tree vertex has
$N$ ordered optional slots numbered $0,\ldots,N-1$; occupying slot $j$ contributes
$q^j$ and attaches a subtree. The root equation is exactly
$w=z\prod_{j=0}^{N-1}(1+q^jw)$. Only finitely many such trees have a fixed number
of vertices. Products of these tree series give the corresponding assertion for
$w^m$.

For $N=2$, ordinary expansion of $(1+t)^n(1+qt)^n$ yields the Narayana formula
\begin{equation}\label{eq:narayana}
 [z^n]w=\frac1n\sum_{j=0}^{n-1}\binom nj\binom n{j+1}q^j.
\end{equation}
The first four coefficients, starting at $z^1$, are
\[
 1,\quad 1+q,\quad 1+3q+q^2,\quad 1+6q+6q^2+q^3.
\]
These are not obtained by replacing the two ordinary binomials in
\eqref{eq:narayana} by Gaussian binomials.

For a reciprocal example, take
$w=z/\poch{aw}qN^r$ with integers $r,N\ge1$. Then
\begin{equation}\label{eq:reciprocal-reversion}
 [z^n]w^m=\frac mn a^{n-m}
 \sum_{\substack{k_1+\cdots+k_{rn}=n-m\\k_i\ge0}}
       \prod_{i=1}^{rn}\qb{N+k_i-1}{k_i}.
\end{equation}
Here \cref{thm:reciprocal} replaces the finite theorem. Formula
\eqref{eq:lagrange-master} continues to apply when $r$ is not an integer, whereas
\eqref{eq:reciprocal-reversion} uses $rn$ separate factors and therefore does not.

\section{A genuinely shifted equation: the area-type \texorpdfstring{$q$}{q}-Catalan series}\label{sec:qcatalan}
The equation in \cref{cor:fuss} contains $q$-dependent coefficients but no dilation
of its unknown series. A true $q$-functional equation behaves differently. The
existence of several inequivalent $q$-Catalan families is well documented;
\cite{Cigler} develops a uniform viewpoint. We give a complete derivation of the
particular normalization needed here.

\begin{theorem}[A quotient solution and a finite coefficient formula]\label{thm:qcat}
Over $\Q(q)$ define
\begin{equation}\label{eq:U-Airy}
 U(z)=\sum_{n\ge0}u_nz^n,\qquad
 u_n=\frac{(-1)^nq^{n(n-1)}}{D_n}.
\end{equation}
Then
\begin{equation}\label{eq:C-Airy}
 C(z)=\frac{U(qz)}{U(z)}
\end{equation}
is the unique formal solution with constant coefficient $1$ of
\begin{equation}\label{eq:area-eqn}
 C(z)=1+zC(z)C(qz).
\end{equation}
For each $n\ge0$, its coefficients have the nonrecursive finite expression
\begin{equation}\label{eq:area-finite}
 [z^n]C(z)=\sum_{k=0}^n q^k u_k
   \sum_{\bm m\in\Pcal(n-k)}
       (-1)^{M(\bm m)}M(\bm m)!
           \prod_{j=1}^{n-k}\frac{u_j^{m_j}}{m_j!},
 \qquad M(\bm m)=\sum_jm_j.
\end{equation}
\end{theorem}
\begin{proof}
The coefficients in \eqref{eq:U-Airy} satisfy, for $n\ge1$,
$(1-q^n)u_n=-q^{2n-2}u_{n-1}$, as follows immediately by cancelling $D_n$.
Consequently $U(z)-U(qz)=-zU(q^2z)$. Division by the unit $U(z)$ gives
\[
 \frac{U(qz)}{U(z)}=1+z\frac{U(q^2z)}{U(qz)}\frac{U(qz)}{U(z)},
\]
which proves \eqref{eq:area-eqn}. Its coefficient of degree $n\ge1$ involves
only lower-degree coefficients on the right, so the solution is unique.
Finally expand $1/U=\sum_{r\ge0}(-1)^r(U-1)^r$ by the ordinary multinomial
theorem and multiply by $U(qz)$. This is exactly \eqref{eq:area-finite}.
\end{proof}

The first terms are
\begin{equation}\label{eq:area-first}
 C(z)=1+z+(1+q)z^2+(1+2q+q^2+q^3)z^3+O(z^4).
\end{equation}
Equation \eqref{eq:area-eqn}, or its binary-tree interpretation, shows that every
coefficient belongs to $\Z_{\ge0}[q]$. Indeed, attach two ordered subtrees to a root
and give the second subtree an extra factor $q$ for each of its vertices. The
resulting weight rule is precisely $zC(z)C(qz)$. This proves that singularities
in the rational summands of \eqref{eq:area-finite} cancel, including at $q=1$.
At $q=1$ the equation is $C=1+zC^2$, so the usual Catalan sequence is recovered.

The polynomials in \eqref{eq:area-first}, the slot family in \eqref{eq:narayana},
and the Gaussian quotient $\qb{2n}n/\qi{n+1}$ are not interchangeable.
For example, their usual Catalan index $n=2$ gives respectively $1+q$,
$1+q$, and $1+q^2$; at $n=3$ the first two already differ. The implicit equation
and normalization, not the name ``$q$-Catalan'', must select the intended family.

\section{A two-variable Gaussian inversion theorem}\label{sec:coupled}

\subsection{A closed coupled example}
\begin{theorem}[Coupled shifted-factorial inverses]\label{thm:coupled}
Let $M,N\ge0$. The system
\begin{equation}\label{eq:coupled-system}
 u=x\poch{-v}qM,\qquad v=y\poch{-u}qN
\end{equation}
has a unique formal solution with $u=x+O((x,y)^2)$ and
$v=y+O((x,y)^2)$. For $m,n\ge1$ and integers $r,s\ge0$ with $m\ge r,n\ge s$,
\begin{equation}\label{eq:coupled-coef}
 [x^my^n]u^rv^s
 =\frac{ms+nr-rs}{mn}
          C_{m-r}^{(n,N)}(q)C_{n-s}^{(m,M)}(q),
\end{equation}
where every $C$ is the finite Gaussian sum in \eqref{eq:Cpn}. The coefficient
is zero if $m<r$ or $n<s$.
On the axes the complete boundary formulas are
\begin{equation}\label{eq:coupled-axes}
 [x^my^0]u^rv^s=\delta_{s0}\delta_{mr},\qquad
 [x^0y^n]u^rv^s=\delta_{r0}\delta_{ns}.
\end{equation}
\end{theorem}
\begin{proof}
Writing $A(t)=\poch{-t}qM$, $B(t)=\poch{-t}qN$, eliminate $v$ to get
$u=xA(yB(u))$. This equation has a unique solution in $x\Q(q)[[y]][[x]]$;
its coefficients define the same bivariate formal series as joint
$(x,y)$-adic iteration. Then $v=yB(u)$.
Ordinary Lagrange--B\"urmann in $x$, with outer function $t^rB(t)^s$, gives
\[
 [x^my^n]u^rv^s
 =\frac1m[t^{m-1}y^{n-s}]
   \bigl(rt^{r-1}B^s+st^rB^{s-1}B'\bigr)A(yB)^m.
\]
Because $[y^{n-s}]A(yB)^m=C_{n-s}^{(m,M)}B^{n-s}$, this becomes
\[
 \frac1m C_{n-s}^{(m,M)}
    [t^{m-r}]\bigl(rB^n+stB^{n-1}B'\bigr).
\]
For $n>0$, differentiation of $B^n$ shows
$[t^{m-r}]tB^{n-1}B'=(m-r)C_{m-r}^{(n,N)}/n$.
The prefactor therefore equals
$(r+s(m-r)/n)/m=(ms+nr-rs)/(mn)$.
Terms with a negative extraction degree are absent. Finally $y=0$ gives
$(u,v)=(x,0)$ and $x=0$ gives $(u,v)=(0,y)$, proving the boundary formulas.
\end{proof}

For $r=1,s=0$, the formula simplifies to
\begin{equation}\label{eq:coupled-u}
 [x^my^n]u=\frac1m C_{m-1}^{(n,N)}(q)C_n^{(m,M)}(q)
 \qquad(m,n\ge1).
\end{equation}
At $q=1$ this is $m^{-1}\binom{Nn}{m-1}\binom{Mm}n$.
With $M=1,N=2$, for instance, $[x^2y]u=1+q$.
The positive integer coefficients also have a two-type rooted-tree interpretation:
u-vertices have $M$ optional slots for v-subtrees, and v-vertices have $N$
optional slots for u-subtrees, weighted by the slot numbers. Thus the rational
prefactor in \eqref{eq:coupled-coef} does not signal nonintegral coefficients.

\subsection{The Good determinant behind the formula}
For a general system $w_i=z_i\Phi_i(\bm w)$ with unit $\Phi_i$, the ordinary
multivariate Lagrange--Good formula reads
\begin{equation}\label{eq:Good}
 [\bm z^{\bm n}]H(\bm w)
 =[\bm t^{\bm n}]H(\bm t)\prod_i\Phi_i(\bm t)^{n_i}
   \det\left(\delta_{ij}-t_j\partial_{t_j}\log\Phi_i(\bm t)\right).
\end{equation}
It follows by the formal change of variables $z_i=t_i/\Phi_i(\bm t)$ in the
multivariate residue: the displayed matrix is exactly the Jacobian on the
logarithmic differentials $dt_j/t_j$. The residue invariance is the
multivariate counterpart of the substitution in \cref{thm:lagrange}, and the
formula is one of the base article's central tools \cite{CCC}.

For \eqref{eq:coupled-system} the determinant is
\[
 1-uv\frac{A'(v)}{A(v)}\frac{B'(u)}{B(u)}.
\]
Extracting the two factors gives
$1-(m-r)(n-s)/(mn)$ times the product of the two $C$ coefficients, which is
\eqref{eq:coupled-coef}. The preceding proof independently derives this special
case without invoking a multivariate residue theorem.

More generally, a factor
$\poch{a\bm t^{\bm\nu}}{q^h}N^\alpha$ has
\begin{equation}\label{eq:Good-entry}
 t_j\partial_{t_j}\log\poch{a\bm t^{\bm\nu}}{q^h}N^\alpha
  =-\alpha\nu_j\sum_{\ell\ge1}
       a^\ell S_N(q^{h\ell})\bm t^{\ell\bm\nu}.
\end{equation}
Here $\bm\nu$ is a nonzero nonnegative integer vector. Thus, after imposing a
fixed multidegree bound, both the determinant and its product factors have
explicit finite expansions: expand the determinant over permutations and expand
the exponential of the logarithmic product over multiplicities indexed by
$(\bm\nu,\ell)$. Only those with $\ell\bm\nu\le\bm n$ can contribute.
This gives a direct route from the source's Good theorem to multivariate
$q$-shifted-factorial identities without changing its Jacobian or introducing
unjustified Jackson derivatives.

\section{Taylor coefficients at \texorpdfstring{$q=1$}{q=1}: Bernoulli--Bell formulas}
\label{sec:one-jets}

Let the ordinary Bernoulli numbers $B_j$ be fixed by
\[
 \frac{t}{e^t-1}=\sum_{j\ge0}B_j\frac{t^j}{j!},\qquad B_1=-\frac12.
\]
These are not the Bell polynomials $B_n(x_1,\ldots,x_n)$; arguments distinguish
our two conventional notations. The Bernoulli normalization agrees with
\cite{DLMF24}. Write
\begin{equation}\label{eq:Delta}
 \Delta_r(n,k)=\sum_{j=1}^n j^r-\sum_{j=1}^k j^r
                  -\sum_{j=1}^{n-k}j^r,\qquad K=k(n-k).
\end{equation}
These are finite sums. In particular $\Delta_1=K$ and
$\Delta_2=K(n+1)$.

\begin{theorem}[All exponential-coordinate derivatives at one]\label{thm:one}
For $0\le k\le n$, define
\begin{equation}\label{eq:kappa-one}
 \kappa_1=\frac K2,\qquad
 \kappa_{2r}=\frac{B_{2r}}{2r}\Delta_{2r}(n,k)\quad(r\ge1),\qquad
 \kappa_{2r+1}=0\quad(r\ge1).
\end{equation}
Then, as a formal Taylor series at $t=0$,
\begin{align}
 \qbb nk{e^t}
 &=\binom nk\exp\left(\sum_{j\ge1}\kappa_j\frac{t^j}{j!}\right),
 \label{eq:one-exp}\\
 \left.\frac{d^r}{dt^r}\qbb nk{e^t}\right|_{t=0}
 &=\binom nk B_r(\kappa_1,\ldots,\kappa_r).
 \label{eq:one-deriv}
\end{align}
Equivalently, its coefficient of $t^r$ is
\begin{equation}\label{eq:one-finite}
 \binom nk
 \sum_{\substack{a+\sum_{j\ge1}2j b_j=r\\a,b_j\ge0}}
 \frac{(K/2)^a}{a!}
 \prod_{j=1}^{\lfloor r/2\rfloor}
       \frac1{b_j!}
       \left(\frac{B_{2j}\Delta_{2j}(n,k)}{2j(2j)!}\right)^{b_j}.
\end{equation}
\end{theorem}
\begin{proof}
Set $A(t)=(e^t-1)/t$, with $A(0)=1$. Differentiating its formal logarithm and
using the Bernoulli generating function gives
\[
 \frac{d}{dt}\log A(t)
  =\frac{e^t}{e^t-1}-\frac1t
  =\frac12+\sum_{r\ge1}\frac{B_{2r}}{(2r)!}t^{2r-1}.
\]
Integration with zero constant gives
\[
 \log A(t)=\frac t2+\sum_{r\ge1}
                 \frac{B_{2r}}{2r(2r)!}t^{2r}.
\]
In the factorial quotient for $\qbb nk{e^t}$, the factors $-jt$ cancel to
$\binom nk$. The remaining logarithm is the sum of $\log A(jt)$ for $1\le j\le n$
minus the two analogous denominator sums. This yields \eqref{eq:kappa-one} and
\eqref{eq:one-exp}. Bell expansion proves \eqref{eq:one-deriv}, and expanding
only its linear and even-degree terms proves \eqref{eq:one-finite}.
\end{proof}

\begin{corollary}[Ordinary derivatives at one]\label{cor:one-ordinary}
Let $s(r,j)$ be the \emph{ordinary signed} first-kind Stirling numbers defined by
$x(x-1)\cdots(x-r+1)=\sum_{j=0}^r s(r,j)x^j$. Then
\begin{equation}\label{eq:ordinary-one-jets}
 \left.\frac{d^r}{dq^r}\qb nk\right|_{q=1}
 =\binom nk\sum_{j=0}^r s(r,j)B_j(\kappa_1,\ldots,\kappa_j).
\end{equation}
The $j=0$ term uses $B_0=1$.
\end{corollary}
\begin{proof}
For $\Theta=q\,d/dq$, one has
$q^r(d/dq)^r=\Theta(\Theta-1)\cdots(\Theta-r+1)
=\sum_js(r,j)\Theta^j$, as is checked on $q^m$ and then by linearity.
At $q=e^t$, $\Theta=d/dt$. Evaluate at $t=0$ and use \eqref{eq:one-deriv}.
\end{proof}
For example,
\begin{align}
 \left.\frac d{dq}\qb nk\right|_{q=1}
 &=\binom nk\frac K2,\label{eq:first-one}\\
 \left.\frac {d^2}{dq^2}\qb nk\right|_{q=1}
 &=\binom nk\left(\frac{K^2}{4}-\frac K2+\frac{K(n+1)}{12}\right).
 \label{eq:second-one}
\end{align}
For $\qb42$, these are $12$ and $22$, in agreement with \eqref{eq:G42}.

There is also a probabilistic interpretation that requires no limit theorem.
Choose a $k$-subset $\{i_1<\cdots<i_k\}$ of $\{0,\ldots,n-1\}$ uniformly and let
$X=\sum i_j-\binom k2$. Its moment generating function is
$\qbb nk{e^t}/\binom nk$, by \eqref{eq:subset}. Thus the $\kappa_j$ are its
cumulants. In particular
\[
 \mathbb E X=\frac{k(n-k)}2,\qquad
 \operatorname{Var}X=\frac{k(n-k)(n+1)}{12},
\]
and every odd cumulant above the first vanishes. This is the cumulant form of
Gaussian reciprocity, since $X$ and $k(n-k)-X$ have the same distribution.

\section{Roots of unity and all local derivatives}\label{sec:roots}

\subsection{Cyclotomic orders and the \texorpdfstring{$q$}{q}-Lucas value}
Let $\Phi_d(q)$ denote the $d$th cyclotomic polynomial.
\begin{proposition}[Exact cyclotomic factorization]\label{prop:cyclotomic}
For $0\le k\le n$,
\begin{equation}\label{eq:cyclotomic}
 \qb nk=\prod_{d=1}^n\Phi_d(q)^{e_d(n,k)},\qquad
 e_d(n,k)=\left\lfloor\frac nd\right\rfloor
       -\left\lfloor\frac kd\right\rfloor
       -\left\lfloor\frac{n-k}d\right\rfloor\in\{0,1\}.
\end{equation}
Thus a Gaussian polynomial has at most a simple zero at any primitive root
of any fixed order. In particular $e_1=0$.
\end{proposition}
\begin{proof}
Factor $1-q^j=-(q^j-1)=-\prod_{d\mid j}\Phi_d(q)$ in the quotient
$D_n/(D_kD_{n-k})$. The signs cancel. Counting multiples of $d$ gives the
exponent, and $n=k+(n-k)$ shows it is zero or one.
\end{proof}

\begin{theorem}[$q$-Lucas evaluation]\label{thm:qlucas}
Let $\zeta$ be a primitive $d$th root of unity, $d\ge1$, and write
$n=ad+b$, $k=cd+r$ with $0\le b,r<d$ and $0\le k\le n$. Then
\begin{equation}\label{eq:qlucas}
 \qbb nk\zeta=\binom ac\qbb br\zeta.
\end{equation}
In particular the value is zero when $r>b$. Equivalently,
$\qb nk\equiv\binom ac\qb br\pmod{\Phi_d(q)}$.
\end{theorem}
\begin{proof}
At $q=\zeta$, the finite product decomposes into full blocks:
\[
 \poch{-z}\zeta n=(1-(-z)^d)^a\poch{-z}\zeta b.
\]
The coefficient of $z^{cd+r}$ on the right is zero if $r>b$; otherwise it is
$(-1)^{c(d+1)}\binom ac\zeta^{\binom r2}\qbb br\zeta$.
On the left it is $\zeta^{\binom k2}\qbb nk\zeta$.
The phases agree because
\[
 \zeta^{\binom{cd+r}2-\binom r2}=(-1)^{c(d+1)}.
\]
For odd $d$, the exponent is a multiple of $d$ and both sides are $1$;
for even $d$, use $\zeta^{d/2}=-1$ and the parity of $c(cd-1)$.
This proves the value identity. Since it holds at every primitive root, the
polynomial difference is divisible by the monic polynomial $\Phi_d$ over $\Z$.
\end{proof}
The congruence is a standard $q$-Lucas theorem; treatments and extensions appear
in \cite{HarmanHopkins,FormichellaStraub}. The proof above includes the phase
factor, which is easy to lose when using the finite theorem at even-order roots.
For example,
\begin{equation}\label{eq:minus-one}
 \qbb nk{-1}=
 \begin{cases}
  0,&n\text{ even and }k\text{ odd},\\
  \displaystyle\binom{\lfloor n/2\rfloor}{\lfloor k/2\rfloor},&\text{otherwise}.
 \end{cases}
\end{equation}

\subsection{Regularized logarithms give every derivative}
The next theorem refines the value formula to a complete local Taylor series.
Only finite sums and products are needed for any prescribed coefficient.
It also covers $d=1$, recovering \cref{thm:one}.

Fix $n,k$ and define signed multiplicities
\begin{equation}\label{eq:epsilon-j}
 \varepsilon_j=\mathbf1_{j\le n}-\mathbf1_{j\le k}
                         -\mathbf1_{j\le n-k}\qquad(1\le j\le n).
\end{equation}
For an integer $\ell\ge1$ put
\begin{equation}\label{eq:gell}
 g_1=\frac12,\qquad g_{2r}=\frac{B_{2r}}{2r}\ (r\ge1),\qquad
 g_{2r+1}=0\ (r\ge1).
\end{equation}
For $a\ne1$ define the finite rational function
\begin{equation}\label{eq:hell}
 h_\ell(a)=-\sum_{v=1}^{\ell}\stir\ell v(v-1)!
                      \left(\frac{a}{1-a}\right)^v,
\end{equation}
where the braces are \emph{ordinary} second-kind Stirling numbers.

\begin{theorem}[All-order root-of-unity expansion]\label{thm:root-jets}
Let $\zeta$ be primitive of order $d$, and let $e=e_d(n,k)\in\{0,1\}$.
Define the nonzero algebraic constant
\begin{equation}\label{eq:root-C}
 C_\zeta=(-1)^e
       \prod_{\substack{1\le j\le n\\d\mid j}}j^{\varepsilon_j}
       \prod_{\substack{1\le j\le n\\d\nmid j}}
                           (1-\zeta^j)^{\varepsilon_j},
\end{equation}
and, for $\ell\ge1$,
\begin{equation}\label{eq:root-kappa}
 K_\ell=\sum_{\substack{1\le j\le n\\d\mid j}}
                  \varepsilon_j j^\ell g_\ell
       +\sum_{\substack{1\le j\le n\\d\nmid j}}
                  \varepsilon_j j^\ell h_\ell(\zeta^j).
\end{equation}
Then in the local coordinate $q=\zeta e^t$,
\begin{equation}\label{eq:root-expansion}
 \qbb nk{\zeta e^t}
   =C_\zeta t^e\exp\left(\sum_{\ell\ge1}K_\ell\frac{t^\ell}{\ell!}\right).
\end{equation}
In particular, if $J_j$ denotes the $j$th $t$-derivative at zero, then
\begin{equation}\label{eq:root-J}
 J_j=\begin{cases}
 0,&j<e,\\
 \displaystyle C_\zeta\frac{j!}{(j-e)!}
                    B_{j-e}(K_1,\ldots,K_{j-e}),&j\ge e.
 \end{cases}
\end{equation}
Every ordinary $q$-derivative is therefore given by the finite sum
\begin{equation}\label{eq:root-ordinary}
 \left.\frac{d^r}{dq^r}\qb nk\right|_{q=\zeta}
       =\zeta^{-r}\sum_{j=0}^r s(r,j)J_j.
\end{equation}
\end{theorem}
\begin{proof}
Use the product $\qb nk=\prod_{j=1}^n(1-q^j)^{\varepsilon_j}$ before taking
any logarithm. When $d\mid j$, write
\[
 1-e^{jt}=-jt\,\frac{e^{jt}-1}{jt}.
\]
The total power of $t$ is $\sum_{d\mid j}\varepsilon_j=e$; the constant factors
produce the first product in \eqref{eq:root-C}, including its sign.
When $d\nmid j$, factor out the nonzero number $1-\zeta^j$.
After removing these constants and $t^e$, the remaining product is a unit with
constant term $1$, so its formal logarithm is legitimate even when $e=1$.

The $\ell$th derivative at zero of
$\log((e^t-1)/t)$ is $g_\ell$, by the proof of \cref{thm:one}.
For the other factors,
\[
 \left.\frac{d^\ell}{dt^\ell}\log(1-ae^t)\right|_{t=0}
   =\left(a\frac d{da}\right)^\ell\log(1-a)=h_\ell(a).
\]
To verify the last equality without an infinite polylogarithm series, use
$(a\partial_a)^\ell=\sum_{v=1}^{\ell}\stir\ell v a^v\partial_a^v$
and $\partial_a^v\log(1-a)=-(v-1)!/(1-a)^v$.
Scaling $t$ by $j$ multiplies the derivative by $j^\ell$.
Summing the signed multiplicities proves \eqref{eq:root-kappa} and
\eqref{eq:root-expansion}. Expanding the exponential by Bell polynomials gives
\eqref{eq:root-J}. Finally use
$q^r\partial_q^r=\sum_js(r,j)(q\partial_q)^j$ at $q=\zeta e^t$ to obtain
\eqref{eq:root-ordinary}.
\end{proof}

If $e=0$, $C_\zeta=\qbb nk\zeta$ and \cref{thm:qlucas} supplies its shorter
value. If $e=1$, the value is zero but
$\left.\partial_q\qb nk\right|_{q=\zeta}=C_\zeta/\zeta\ne0$.
For instance $\qb41=1+q+q^2+q^3$ has, at $\zeta=-1$, $e=1$ and
$C_{-1}=-2$, giving derivative $2$.
For $d=1$, all factors are in the divisible class, $e=0$,
$C_1=\binom nk$, and $K_\ell$ equals \eqref{eq:kappa-one}.

This theorem illustrates the full coefficient-calculus chain: cyclotomic
factorization isolates the zero, Bernoulli and second-kind Stirling numbers
compute the regular logarithmic derivatives, Bell polynomials exponentiate them,
and signed first-kind Stirling numbers change from the logarithmic coordinate
back to ordinary derivatives. Direct substitution into $D_n/(D_kD_{n-k})$ would
instead produce uninformative $0/0$ expressions.

\section{Exact geometric extrapolation and mode annihilation}\label{sec:extrapolation}

\begin{theorem}[Gaussian extrapolation weights]\label{thm:extrapolation}
Let $N\ge0$, with $q\ne0$ and $q^j\ne1$ for $1\le j\le N$. Define
\begin{equation}\label{eq:weights}
 w_j^{(N)}(q)=
   \frac{(-1)^{N-j}q^{\binom{N-j+1}2}}{D_jD_{N-j}}
 =\frac{(-1)^{N-j}q^{\binom{N-j+1}2}}{D_N}\qb Nj,
 \qquad 0\le j\le N.
\end{equation}
Then
\begin{equation}\label{eq:weight-polynomial}
 \sum_{j=0}^Nw_j^{(N)}(q)t^j
       =\frac{\prod_{r=1}^N(t-q^r)}{D_N},
\end{equation}
and
\begin{equation}\label{eq:annihilation}
 \sum_{j=0}^Nw_j^{(N)}(q)q^{j\ell}
   =\begin{cases}1,&\ell=0,\\0,&1\le\ell\le N.\end{cases}
\end{equation}
Consequently, if $a_j=L+\sum_{\ell=1}^N c_\ell q^{j\ell}$, then
$L=\sum_{j=0}^Nw_j^{(N)}a_j$ exactly.
\end{theorem}
\begin{proof}
The coefficient of $t^j$ in $\prod_{r=1}^N(t-q^r)$ is
$(-1)^{N-j}$ times the elementary symmetric polynomial of degree $N-j$ on
$q,q^2,\ldots,q^N$. By \cref{thm:finite}, this is
$(-1)^{N-j}q^{\binom{N-j+1}2}\qb Nj$.
Division by $D_N$ gives \eqref{eq:weight-polynomial}.
Set $t=1$ to obtain weight sum $1$, and set $t=q^\ell$ to obtain the zero moments.
The final statement follows by linearity.
\end{proof}

Equivalently, these are the Lagrange interpolation weights for evaluating a
polynomial at zero from its values at $1,q,\ldots,q^N$. Thus any polynomial
$p$ of degree at most $N$ satisfies
$p(0)=\sum_jw_j^{(N)}p(hq^j)$ for any scalar $h$.
In particular, for dyadic refinement $q=1/2$ and $N=2$,
\begin{equation}\label{eq:dyadic-weights}
 (w_0,w_1,w_2)=\left(\frac13,-2,\frac83\right),\qquad
 L=\frac13a_0-2a_1+\frac83a_2
\end{equation}
whenever $a_j=L+c_1 2^{-j}+c_2 4^{-j}$.
This is an exact finite-spectrum statement, not an assertion that every
convergent approximation sequence admits exact extraction from three entries.

For an analytic $f(h)=\sum_{\ell\ge0}c_\ell h^\ell$ and arguments inside its disk
of convergence, the exact remainder is
\begin{equation}\label{eq:extrap-error}
 \sum_{j=0}^Nw_j^{(N)}f(hq^j)-f(0)
   =\sum_{\ell\ge N+1}c_\ell h^\ell
       \frac{\prod_{r=1}^N(q^\ell-q^r)}{D_N}.
\end{equation}
This follows by applying \eqref{eq:weight-polynomial} termwise to the finite
weighted sum. It implies an $O(h^{N+1})$ error for fixed $N,q$. It does not claim
uniformity as $N$ grows or $q$ approaches a forbidden root.

\section{Formal identities, analytic domains, and specialization}\label{sec:analytic}

\subsection{What convergence is, and is not, required}
Finite Gaussian identities, triangular transforms, and coefficient extraction in
$R[[z]]$ need no analytic convergence. Formal exponentials, logarithms, and divided
Lagrange formulas do require the indicated rational denominators. For products
with infinitely many geometric factors, two valid settings have been used:
coefficientwise $q$-adic products, or rational logarithmic coefficients such as
$S_\infty(q^j)=1/(1-q^j)$. Neither should be mislabeled as ordinary $z$-adic
convergence of the finite products.

For a fixed complex $|q|<1$, $\poch{az}q\infty$ converges locally uniformly in
$z$: on each compact set, $\sum_{j\ge0}|azq^j|$ converges uniformly. Away from its
zeros, its reciprocal is analytic. Fractional powers of a unit product have a
unique analytic branch near $z=0$ with value $1$, agreeing with the formal power
used here. The logarithm computation is justified in a sufficiently small disk
by absolute convergence, and the identities elsewhere follow by analytic
continuation where both sides are defined. These conditions recover the standard
analytic $q$-binomial theorem \cite{DLMF17,DLMF175}.

If $\Phi$ in \cref{thm:lagrange} is analytic near zero and $\Phi(0)=1$, the
analytic inverse-function theorem makes its formal solution an analytic germ.
The displayed coefficient formulas are then its actual Taylor coefficients.
No claim is made that they converge past the first singularity of that germ.
The two-variable system has the analogous local analytic interpretation because
its implicit Jacobian at the origin is the identity.

\subsection{A specialization checklist}
Polynomial identities such as \eqref{eq:finite}, \eqref{eq:vandermonde},
\eqref{eq:transform-inverse}, and \eqref{eq:qlucas} survive specialization at any
$q$. By contrast, three types of rational formulas need attention.
First, interpolation on geometric or $q$-integer nodes requires distinct nodes;
colliding nodes require a confluent limit rather than raw substitution.
Second, $q$-divided-power coordinates need invertible $\qf j$, so a coordinate
formula can have genuine poles even when a separately defined polynomial family
does not. Third, infinite-factor logarithms involve $1-q^j$ in denominators,
and a root-of-unity specialization is not available without a separate limiting
argument or a proved cancellation in the final coefficients.

The safe procedure for a Gaussian polynomial is to use its subset expansion,
its cyclotomic factorization, or \cref{thm:root-jets}. The fact that an intermediate
formula has $0/0$ terms is not evidence that the Gaussian polynomial is undefined.
Conversely, polynomiality of a Gaussian coefficient does not regularize every
expression built from reciprocals of $q$-factorials.

\section{Implementation, verification, and integration with the source}\label{sec:verification}

\subsection{Finite algorithms suggested by the formulas}
The partition kernel \eqref{eq:E-kernel} uses $p(n)$ multiplicity profiles, where
$p(n)$ is the ordinary partition number. It is often much smaller than expanding
all ordered compositions, but it is not claimed to be the fastest way to compute
long coefficient ranges. Product multiplication and triangular recurrences may
be more efficient for bulk computation. An explicit nonrecursive formula and an
asymptotically optimal algorithm are different objectives.

For one coefficient, a dependable implementation follows this sequence:
construct the finitely many logarithmic inputs needed at that degree; enumerate
integer partitions; evaluate the products in exact rational or symbolic
arithmetic; and simplify only the completed expression before specializing a
potentially singular $q$. For \cref{thm:root-jets}, first separate the divisible
and nondivisible factors, then form $K_1,\ldots,K_r$. No infinite series needs to
be stored to obtain a derivative of order $r$.

The accompanying \texttt{verify\_identities.py} implements independent checks of
the displayed identities. Gaussian polynomials are constructed from subset
weights rather than from the identities being tested. Product expansions are
compared with partition kernels; fixed-point series are compared with Lagrange
and coupled formulas; and polynomial derivatives are compared with the
Bernoulli--Bell and regularized root formulas. The archive includes the actual
machine-readable test report. The included run passed 2,693 exact assertions in four groups. The finite
ranges, exact coefficient fields, and per-identity counts are recorded there;
these are finite regression tests, not an unlimited formal verification.

\Needspace{12\baselineskip}
\subsection{Normalization audit}
\begin{longtable}{@{}>{\RaggedRight\arraybackslash}p{0.25\textwidth}>{\RaggedRight\arraybackslash}p{0.34\textwidth}>{\RaggedRight\arraybackslash}p{0.34\textwidth}@{}}
\toprule
\textbf{Operation} & \textbf{Correct object} & \textbf{Incorrect shortcut}\\
\midrule
\endfirsthead
\toprule
\textbf{Operation} & \textbf{Correct object} & \textbf{Incorrect shortcut}\\
\midrule
\endhead
Finite product & $q^{\binom k2}\qb Nk$ & Omitting $q^{\binom k2}$\\[1mm]
Gaussian inverse & $(-1)^r q^{\binom r2}\qb nr$ & Replacing the inverse by $T_{-c}$\\[1mm]
Ordinary composition & Ordinary $k!/\prod m_j!$ inside \eqref{eq:Bq-def} & Replacing every profile multiplicity by a Gaussian multinomial\\[1mm]
Lagrange--B\"urmann & $m/n$ and an ordinary derivative & Replacing $n$ by $\qi n$ without a new theorem\\[1mm]
Newton limit & $q$-integer nodes tend to integers & Treating geometric nodes as if they did not coalesce\\[1mm]
Root-of-unity values & Polynomial evaluation or regularization & Substituting into an uncancelled factorial quotient\\[1mm]
$q$-Catalan family & Its defining equation and normalization & Assuming all Gaussian and area refinements coincide\\[1mm]
Infinite product & $q$-adic, logarithmic, or analytic interpretation & Claiming stabilization in the $z$-adic topology over $\Q(q)$\\
\bottomrule
\end{longtable}

\subsection{Where the extension attaches}
The following crosswalk uses labels observed in the canonical source, not its
mutable page or theorem numbers. Labels are pointers to mathematical interfaces,
not assertions that the new theorems have existing Lean counterparts.
\begin{longtable}{@{}>{\RaggedRight\arraybackslash}p{0.36\textwidth}>{\RaggedRight\arraybackslash}p{0.57\textwidth}@{}}
\toprule
\textbf{Base-source interface} & \textbf{Extension developed here}\\
\midrule
\endfirsthead
\toprule
\textbf{Base-source interface} & \textbf{Extension developed here}\\
\midrule
\endhead
\texttt{thm:newton-\allowbreak{}expansion} & Two node systems, explicit Gaussian interpolation, and the specified $q$-Stirling inverse arrays (\cref{sec:newton}).\\[1mm]
\texttt{thm:bell-\allowbreak{}symmetric-\allowbreak{}functions} & Geometric-alphabet specialization and the ramified shifted-factorial Bell kernel (\cref{sec:master}).\\[1mm]
\texttt{thm:merged-\allowbreak{}binomial-\allowbreak{}inversion} & Gaussian inverse weights and subspace-lattice M\"obius interpretation (\cref{sec:transforms}).\\[1mm]
\texttt{thm:merged-\allowbreak{}weighted-\allowbreak{}binomial-\allowbreak{}translation} & The corrected composition law; the ordinary additive law does not persist (\cref{sec:transforms}).\\[1mm]
\texttt{eq:merged-\allowbreak{}bell-\allowbreak{}normalization} & Coordinate-specific Bell conversion, with ordinary profile multiplicities retained (\cref{sec:composition}).\\[1mm]
\texttt{thm:lagrange-\allowbreak{}burmann} & Recurrence-free product-weighted inverse coefficients and coupled inverses (\cref{sec:lagrange,sec:coupled}).\\
\bottomrule
\end{longtable}

\subsection{Mathematical and formalization status}
The classical finite theorem, Vandermonde and Chu--Vandermonde sums,
Gaussian inversion, and $q$-Lucas evaluation are not presented as discoveries.
The master kernel and its applications provide a single reproducible route to
many further finite formulas. The root-of-unity derivative theorem is proved
here in the precise normalization stated, without claiming historical novelty
for its regularization method. All proofs are ordinary mathematical proofs.
No Lean files were added or compiled for this extension. Exact finite tests are
regression checks, not substitutes for universally quantified proofs.

The formulas also identify a natural formalization order: first polynomial
Gaussian identities and denominator-free transforms, then finite partition
kernels, then normalized formal-series composition and Lagrange inversion,
and finally cyclotomic regularization. This order keeps algebraic prerequisites
separate from analytic convergence and from singular coordinate changes.

\section{Conclusion}
Gaussian binomial coefficients are not merely ornamented binomial coefficients.
They are coefficient extractors for geometric alphabets, structure constants for
$q$-factorial convolution, and evaluation matrices for $q$-integer Newton bases.
These roles connect directly with ordinary Bell polynomials, ordinary Stirling
numbers, and ordinary Lagrange inversion once the coefficient coordinates are
stated correctly.

The central computational principle is
\[
 \text{geometric product}
 \ \longrightarrow\ \text{finite logarithmic inputs}
 \ \longrightarrow\ \text{Bell partition sum}.
\]
For inversion, the same inputs are multiplied by the Lagrange index; for
multivariate inversion, they are supplemented by a finite Jacobian determinant.
At roots of unity, a preliminary cyclotomic factorization removes the zero before
the Bell mechanism is used. This yields explicit formulas for coefficients,
compositions, inverse series, and local derivatives while making clear which
identities extend polynomially to singular parameter values and which do not.

\begin{thebibliography}{99}
\bibitem{CCC}
Vladimir Reshetnikov, ProveIt repository,
\emph{Combinatorial Coefficient Calculus: Stirling, Bell, Eulerian, and
Bernoulli--Euler Families; Bell-Polynomial and Fa\`a di Bruno Calculus;
Lagrange--B\"urmann and Multivariate Good Inversion}, canonical \TeX{} source.
Consulted September 4, 2026.
\url{https://github.com/VladimirReshetnikov/ProveIt/tree/main/Analysis/FabiusFunction/docs/semi-formalized-research-frontiers/drafts/combinatorial-coefficient-calculus/Combinatorial_Coefficient_Calculus}.

\bibitem{provenance}
ProveIt repository, \emph{Provenance and consolidation boundary},
\texttt{PROVENANCE.md} in the canonical package cited above.
Consulted September 4, 2026. This distinguishes the evolving source from retained
historical PDF renders and from result-specific formalization coverage.

\bibitem{DLMF17}
NIST Digital Library of Mathematical Functions, \S17.2,
\emph{Calculus}, especially the $q$-shifted factorials, $q$-binomial coefficients,
finite and infinite binomial theorems, and Jackson calculus.
\url{https://dlmf.nist.gov/17.2}. Consulted September 4, 2026.

\bibitem{DLMF175}
NIST Digital Library of Mathematical Functions, \S17.5,
\emph{$ {}_0\phi_0,{}_1\phi_0,{}_1\phi_1$ Functions}.
\url{https://dlmf.nist.gov/17.5}. Consulted September 4, 2026.

\bibitem{DLMF176}
NIST Digital Library of Mathematical Functions, \S17.6,
\emph{$ {}_2\phi_1$ Function}, terminating $q$-Chu--Vandermonde sums.
\url{https://dlmf.nist.gov/17.6}. Consulted September 4, 2026.

\bibitem{DLMF26}
NIST Digital Library of Mathematical Functions, \S26.9,
\emph{Integer Partitions: Restricted Number and Part Size}.
\url{https://dlmf.nist.gov/26.9}. Consulted September 4, 2026.

\bibitem{DLMF24}
NIST Digital Library of Mathematical Functions, \S24.2,
\emph{Definitions and Generating Functions} for Bernoulli numbers and polynomials.
\url{https://dlmf.nist.gov/24.2}. Consulted September 4, 2026.

\bibitem{FormichellaStraub}
Sam Formichella and Armin Straub,
\emph{Gaussian binomial coefficients with negative arguments},
arXiv:1802.02684 (2018).
\url{https://arxiv.org/abs/1802.02684}.

\bibitem{HarmanHopkins}
Nate Harman and Sam Hopkins,
\emph{Quantum integer-valued polynomials},
arXiv:1601.06110 (2016).
\url{https://arxiv.org/abs/1601.06110}.

\bibitem{Cigler}
Johann Cigler,
\emph{$q$-Catalan numbers and $q$-Narayana polynomials},
arXiv:math/0507225 (2005).
\url{https://arxiv.org/abs/math/0507225}.
\end{thebibliography}
\end{document}
